\section{Supporting Proofs}
\label{app:proofs}
\subsection{Rounded-noise tails}
\label{app:lem-rounded-tail}
\begin{proof}
The rounding cells for \(-h,\ldots,h\) fill
\([-h-1/2,h+1/2)\), apart from boundary conventions. Hence
\(|Z|>h\) iff \(|Y|\ge h+1/2\), up to half-integer ties,
which have probability zero under the continuous Gaussian.
Integrating its two symmetric tails proves \eqref{eq:rounded-strict-tail}.
The non-strict identity follows from integer-valuedness:
\(|Z|\ge h\) iff \(|Z|>h-1\).
\end{proof}

\subsection{Block-support noise bound}
\label{app:lem-block-noise}
\begin{proof}
There are \(m_j\) valid segment positions in block \(j\), each in
\([0,B-1]\), and all other coefficients are zero. Since
\((L-1)n<J\le Ln\), actual blocks are nonempty. Unused high bits in the
last segment only reduce the stated bound.
For degree-\(<n\) representatives, the coefficient of a negacyclic product is
\[
 [X^k](Pe)=\sum_{u+v=k}P_u e_v-\sum_{u+v=k+n}P_u e_v.
\]
The minus sign is exactly the reduction \(X^n=-1\).
For each \(u\) there is one \(v=(k-u)\bmod n\) in these two sums.
The triangle inequality therefore bounds its absolute value by
\(\sum_u|P_u|\|e\|_\infty\).
On \(G_h\), each \(\|e_i\|_\infty\le h\); summing
\(\|\widehat P_{i,j}e_i\|_\infty\le m_j(B-1)h\) over all \(N\) records
proves \eqref{eq:aggregate-good-event-bound}.
\end{proof}

\subsection{Normal-form specialization}
\label{app:normal-form}
\begin{proof}[Specialization sketch]
Given \((a_i,b_i)_{i=0}^{\tau}\), output zero if \(a_0\) is not a unit.
Otherwise give \(\mathcal D\) the pairs
\[
 A_i=-a_0^{-1}a_i,\qquad B_i=b_i+A_i b_0,\quad i\in[\tau].
\]
In the real case, \(B_i=A_i e_0+e_i\), so the new shared secret is
\(e_0\leftarrow\chi\) and all remaining errors still have law \(\chi\).
The unit event depends only on \(a_0\), hence does not bias \(e_0\);
unit multiplication preserves the independent uniform \(A_i\).
In the uniform case, conditional on \(a_0,b_0\) and the unit event,
the transformed pairs are independent uniform pairs.
Thus \(\mathcal B\)'s gap is exactly \(p_{\rm unit}\) times
\(\mathcal D\)'s gap. Unit testing and inversion are polynomial-time
operations: use the multiplication matrix of \(a_0\) over \(\mathbb Z_q\),
test whether its determinant is coprime to \(q\), and invert when it is.
No factorization of \(q\) is required.
\end{proof}

\subsection{Ciphertext scaling and translation}
\label{app:ske}
\begin{proof}
For each normal-form pair \((A,B=As+E)\) and canonical message lift
\(\widehat m\), set
\[
 (a,b)=(tA,tB+\widehat m).
\]
Then \(b=as+tE+\widehat m\), exactly XPIR's SKE syntax, or
\(C=(b,-a)\) in our ordering. In the uniform branch,
\((tA,tU+\widehat m)\) remains uniform because \(t\) is a unit modulo \(q\).
Apply the same bijection to all samples sharing \(s\).
This is the noise-scaling observation of
\cite[Proposition~1]{homomorphic}, followed by message translation;
the argument uses only \(\gcd(t,q)=1\), including for composite \(q\).
\end{proof}

\subsection{Single-challenge privacy reduction}
\label{app:privacy}
\begin{proof}
First construct a normal-form distinguisher \(\mathcal D\).
Given \(pp\) and \(N\) normal-form-or-uniform pairs \((A_i,B_i)\),
it runs \(\mathcal A_1(pp)\) to obtain \((M,I_0,I_1,\sigma)\). If the inputs are
illegal, it outputs a fair bit, matching the experiment's convention.
Otherwise it privately chooses \(\beta\leftarrow\{0,1\}\) and forms
\begin{equation}
\label{eq:reduction-query}
 C_i=(tB_i+v_{I_\beta,i},-tA_i),\qquad Q=(C_1,\ldots,C_N).
\end{equation}
The constants \(v_{I_\beta,i}\) use their canonical message lifts.
It computes \(P=\mathsf{DBEncode}(M)\) and
\(\mathsf{Resp}=\mathsf{ReplyGen}(pp,P,Q)\), gives
\((Q,\mathsf{Resp})\) to \(\mathcal A_2\) with precisely the other inputs
of the experiment, and outputs one iff the guess \(\beta'\) equals \(\beta\).
It needs neither \(s\) nor an encryption or decryption oracle.

For real normal-form samples, \(B_i=A_i s+E_i\), with one shared
\(s\leftarrow\chi\) and fresh independent \(E_i\leftarrow\chi\).
By Corollary~\ref{cor:ske-pseudorandom}, setting \(a_i=tA_i\) gives
\[
 C_i=(a_i s+tE_i+v_{I_\beta,i},-a_i),
\]
exactly a fresh \(\mathsf{Enc}_s(v_{I_\beta,i})\). This preserves the
entire joint law of all \(N\) query ciphertexts, including their shared
secret, not just individual marginals. The first-stage output has the
same law as in the game because it depends only on \(pp\) and adversary
randomness, independently of the samples and of \(\beta\).
The encoded database and reply use the actual deterministic protocol
algorithms. Hence \((Q,\mathsf{Resp})\) has exactly the formal protocol
view's distribution, and
\[
 \Pr[\mathcal D(pp,\mathcal F_{\rm real}^{(N)})=1]
   =\Pr[\text{privacy experiment succeeds}].
\]

For uniform samples, \(A_i,B_i\) are independent uniform ring elements.
Multiplication by the unit \(t\), negation, and translation by the fixed
\(v_{I_\beta,i}\) preserve uniformity. Conditional on every legal first-stage
output and either value of \(\beta\), the entire query is uniform in
\((R_q^2)^N\), independently of \(\beta\). The reply is deterministic PPT
post-processing of this query and the fixed database. Therefore
\((Q,\mathsf{Resp})\) remains independent of \(\beta\).
The same is true of the adversary's resulting guess. Accounting also
for the independent fair output on illegal input,
\[
 \Pr[\mathcal D(pp,\mathcal F_{\rm unif}^{(N)})=1]=\tfrac12.
\]
Under the success-minus-\(1/2\) convention of the privacy game and the
acceptance-gap convention for PLWE, subtraction yields
\[
 \Adv^{\rm priv}_{\Pi,\mathcal A}(\kappa)
 =\Adv^{\rm NF}_{\mathcal D}(\kappa;N).
\]
Compose \(\mathcal D\) with Corollary~\ref{cor:normal-form} at \(\tau=N\)
to obtain \(\mathcal B\). It consumes \(N+1\) uniform-secret ring samples,
uses the first for normal form, and passes the remaining \(N\) transformed
pairs to \(\mathcal D\). Conditional on the unit event its simulation is
the one just analyzed. With output zero on abort, its uniform-branch
acceptance probability is \(p_{\rm unit}/2\), not \(1/2\); both branches
are scaled by the same \(p_{\rm unit}\). Thus
\[
 \Adv^{\rm uPLWE}_{\mathcal B}(\kappa;N+1)
 =p_{\rm unit}\,\Adv^{\rm NF}_{\mathcal D}(\kappa;N),
\]
which gives the stated bound. There is no extra factor \(2\), \(N\), or
\(2\alpha\): the reduction simulates the joint query in one step.

All steps are polynomial-time: the normal-form operations are as in
Corollary~\ref{cor:normal-form}, both adversary stages are PPT, and
encoding and reply generation operate on a polynomial-size database
and query using polynomially many ring operations with polynomial
bit complexity. Deterministic reply generation adds no distinguishing loss.
\end{proof}
